\documentclass[10pt,a4paper]{amsart}
\usepackage[margin=19mm]{geometry}
\usepackage{amsmath,amssymb,mathtools}
\usepackage[scr=rsfs,cal=euler]{mathalfa}
\usepackage{xfrac}
\usepackage[unicode,hidelinks,bookmarks=false]{hyperref}
\newcommand{\ie}{i.e.\ }
\newcommand{\cf}{cf.\ }
\newcommand{\resp}{respectively\ }
\newcommand{\Subsection}{\subsection}
\providecommand{\nobreakdash}{\nobreak\hbox{-}}
\setlength{\emergencystretch}{2em}
\multlinegap0pt
\usepackage{bm}
\usepackage{bbm}
\usepackage{dsfont}
\usepackage{xspace}
\usepackage{cancel}
\usepackage{mathtools}
\usepackage{amsthm}
\makeatletter
\@ifundefined{thm}{\newtheorem{thm}{Thm}}{}
\makeatother
\makeatletter
\expandafter\ifx\csname clmproof\endcsname\relax
\newenvironment{clmproof}{\begin{proof}[Proof of Claim:]\renewcommand{\qedsymbol}{\rule{1ex}{1ex}}}{\end{proof}}
\fi
\makeatother
\title[3-packings in Triangulations: Algorithms, bounds, and Comple]{3-packings in Triangulations: Algorithms, bounds, and Complexity}
\author{Prosenjit Bose, Anil Maheshwari, Bobby Miraftab, Yota Otachi}
\date{Source: arXiv:2606.29743v1 (CC BY 4.0)}
\begin{document}
\maketitle

\noindent\emph{Source and licence.} 3-packings in Triangulations: Algorithms, bounds, and Complexity. Version arXiv:2606.29743v1 (CC BY 4.0), distributed under the Creative Commons Attribution 4.0 International licence, as recorded in the arXiv licence metadata for this exact version and shown on the abstract page https://arxiv.org/abs/2606.29743v1. Full rights notice: licenses/arxiv-2606.29743-CC-BY-4.0.txt.\par\smallskip

\noindent\textit{Editorial scope.} Counted unit: the Thm whose source label is \texttt{thm:4-conn} in the article (source0.tex lines 1250--1272, source label \texttt{thm:4-conn}), delivered together with the article's own complete original proof of it (source0.tex lines 1274--1327, one \texttt{proof} environment of 54 lines). No mathematics is written, repaired or re-derived: statement and proof are the article's own text line for line. Required context retained uncounted: none. Every definition, display and label of those lines is retained unchanged. Omitted from this package and not redistributed: 1--1249, 1273--1273, 1328--1890. Nothing inside a retained excerpt is omitted or altered: every retained body is delivered exactly as the article's lines are written, so no omission receipt of any kind accompanies this package. Citations: the retained excerpts carry no citation command at all, so no bibliography entry of the article is transcribed and no reference is redistributed. Reference adaptation: none was needed and none was made; every reference inside the delivered unit resolves inside it, so no unresolved reference or citation is produced. Printed slips kept literal: no printed word, symbol, hypothesis or number of the counted statement or of its proof was altered. Layout and class: the article's own class is \texttt{patmorin} with packages including pat, fontenc, inputenc, todonotes, float, subcaption, algorithm, algpseudocode, lineno, natbib, enumitem; the wrapper is the lane's pinned \texttt{amsart} editorial wrapper at 10pt on A4, which restates the theorem-like environments the retained text needs and adds the article's own macro definitions that the retained text needs (none), with extra wrapper packages (bm, bbm, dsfont, xspace, cancel, mathtools). The delivered numbering is the unit's own. Assets: no retained span contains a figure environment, a graphics-inclusion command, a PDF-page inclusion, a table or a TikZ picture; no image file is copied into this package, the package declares no asset dependency and no image file is redistributed with it. Licence: version arXiv:2606.29743v1 of this article is distributed under the Creative Commons Attribution 4.0 International licence, as recorded in the arXiv licence metadata for this exact version and shown on the abstract page https://arxiv.org/abs/2606.29743v1. A full rights notice is delivered at licenses/arxiv-2606.29743-CC-BY-4.0.txt. Characters: every retained excerpt is pure ASCII, so the delivered engine, XeLaTeX with its default Latin Modern text font, reports no missing character.
\begin{thm}\label{thm:4-conn}
Let $G$ be a $4$-connected plane triangulation.
Then the following are equivalent.

\begin{enumerate}
\item $G$ has a triangle factor.

\item $G^*$ has an independent set
$X\subseteq V(G^*)$ such that every face of $G^*$ contains exactly one vertex
of $X$.

\item There exist sets of bounded faces
$\mathcal T^+\subseteq F_b(G^+)$ and $\mathcal T^-\subseteq F_b(G^-)$ such that
every vertex $v_j$ of $C$ belongs to exactly one face of
$\mathcal T^+\cup\mathcal T^-$.

\item There exist sets $X^+\subseteq V(T^+)$ and $X^-\subseteq V(T^-)$ such
that, for every $j\in\{1,\dots,n\}$,
\[
|X^+\cap V(P_j^+)|+|X^-\cap V(P_j^-)|=1.
\]
\end{enumerate}
\end{thm}

\begin{proof}
Since $G$ is $4$-connected, every triangle of $G$ is facial. Hence a
triangle factor of $G$ is exactly a set of pairwise vertex-disjoint faces of
$G$ covering $V(G)$.

$(1)\Rightarrow(2)$.
Let $\mathcal T$ be a triangle factor of $G$, and let
$X=\{f^*:f\in\mathcal T\}\subseteq V(G^*)$.
If two vertices of $X$ were adjacent in $G^*$, then the corresponding faces
of $G$ would share an edge, hence two vertices, contradicting the
vertex-disjointness of $\mathcal T$. So $X$ is independent.
For each $v\in V(G)$, the faces of $G$ incident with $v$ correspond to
the vertices on the dual face of $G^*$ associated with $v$. Since
$\mathcal T$ covers $v$ exactly once, so the dual face contains exactly one
vertex of $X$. Thus $X$ is an independent set.

$(2)\Rightarrow(1)$.
Let $X\subseteq V(G^*)$ be independent and let $\mathcal T=\{f:f^*\in X\}$.
If $f,g\in\mathcal T$ shared a vertex $v$, then both $f^*$ and $g^*$
would lie on the dual face corresponding to $v$, contradicting interdependence. 
Hence the faces in $\mathcal T$ are pairwise vertex-disjoint.
Again by face-exactness, each $v\in V(G)$ is incident with exactly one face
of $\mathcal T$. Therefore $\mathcal T$ is a triangle factor.

$(2)\Leftrightarrow(3)$.
The bounded faces of $G^+$ together with the bounded faces of $G^-$ are
precisely the faces of the triangulation $G$; the outer faces of $G^+$ and
$G^-$ are not faces of $G$ and are not included. Therefore a set
$X\subseteq V(G^*)$ is the same thing as a pair of sets
$(\mathcal T^+,\mathcal T^-)$ with
\[
\mathcal T^\sigma\subseteq F_b(G^\sigma).
\]
For a fixed $j$, the dual face of $G^*$ corresponding to $v_j$ consists
exactly of the faces of $G$ incident with $v_j$, equivalently the bounded
faces of $G^+$ and $G^-$ incident with $v_j$.
Hence that dual face contains exactly one vertex of $X$ if and only if
$v_j$ belongs to exactly one face of $\mathcal T^+\cup\mathcal T^-$.
So (2) and (3) are equivalent.

$(3)\Leftrightarrow(4)$.
For each $\sigma\in\{+,-\}$, vertices of $T^\sigma$ are bounded faces of
$G^\sigma$. Thus a set of bounded faces
$\mathcal T^\sigma\subseteq F_b(G^\sigma)$ corresponds uniquely to a vertex
set $X^\sigma\subseteq V(T^\sigma)$.
By definition of $P_j^\sigma$, a bounded face of $G^\sigma$ contains $v_j$
if and only if the corresponding vertex of $T^\sigma$ lies on $P_j^\sigma$.
Therefore $v_j$ belongs to exactly one face of
$\mathcal T^+\cup\mathcal T^-$ if and only if
\[
|X^+\cap V(P_j^+)|+|X^-\cap V(P_j^-)|=1.
\]
So (3) and (4) are equivalent.
\end{proof}

\end{document}
